% ======================================================================
% capitulos/linea2/capitulo6.tex
% ----------------------------------------------------------------------
% LÍNEA CE02 — INGENIERÍA DE SOFTWARE
% Evidencia CE0211: Especificación de requerimientos validada
% Fuente: WISP_Entregable1_Linea3.docx + verificación contra el código real
%         (repositorio Wisp, rama main, 18/11/2026)
% ======================================================================
\section{Sección 1: Especificación de Requerimientos}
\label{chap:ce02-requerimientos}

Esta sección corresponde a la
evidencia \textbf{CE0211 — Especificación de requerimientos validada}. Define qué
debe hacer el \textbf{Sistema de Gestión Integral para ISP de Fibra Óptica
Multi-Zona} (en adelante, \emph{el sistema}) a partir del modelado de procesos TO-BE
(\cref{chap:procesos}) y de la solución TIC (\cref{chap:soluciontic}) de la línea
CE01.

\begin{upeubox}{Fundamento metodológico}
La especificación sigue la estructura de la norma ISO/IEC/IEEE~29148:2018
(\emph{Requirements engineering}): requerimientos funcionales agrupados por módulo,
atributos de calidad (no funcionales), reglas de negocio, restricciones e historias
de usuario con criterios de aceptación verificables. Cada requerimiento tiene un
identificador único que se rastrea hasta la arquitectura, el diseño y la base de
datos (\cref{sec:ce02-trazabilidad}).
\end{upeubox}

\begin{upeubox}{Especificación contrastada con el producto real}
A diferencia de un SRS escrito solo en papel, esta versión se \textbf{contrastó con el
código fuente} del repositorio del sistema (rama \texttt{main} al 18/11/2026). Cada
requerimiento indica su estado: \textbf{Impl.} (implementado, con evidencia en el
código), \textbf{Parcial} (implementado en parte o de otra forma) o
\textbf{Pend.} (pendiente, planificado para el semestre~2). Así el documento no
declara como hecho algo que el producto todavía no hace.
\end{upeubox}

% ======================================================================
\subsection{Contexto, problema y usuarios}
\label{sec:ce02-contexto}

WISP es un proveedor de internet por fibra óptica con sede en Juliaca (Puno) que
atiende a sus clientes en 13 zonas rurales y periurbanas, cada una con su propio
router MikroTik RB4011 y una OLT GPON. Antes del sistema, el proceso completo
—captación, instalación, inventario de equipos, cobranza, cortes y comunicación con
el cliente— se gestionaba de forma manual y dispersa (cuadernos, hojas de cálculo,
grupos de WhatsApp y acceso directo a cada router). Esto provocaba errores de
coordinación entre ventas y técnicos, pérdida de trazabilidad de qué borne y qué
equipo quedó en cada domicilio, retrasos en la cobranza y cortes inconsistentes.

\begin{table}[htbp]
\centering
\caption{Usuarios del sistema (roles reales implementados).}
\label{tab:ce02-roles}
\small
\begin{tabularx}{\textwidth}{L{3.2cm}L{3cm}X}
\toprule
\textbf{Rol} & \textbf{Canal} & \textbf{Responsabilidad principal} \\
\midrule
Administrador (\texttt{admin}) & Panel web & Gerente/ingeniero: valida ventas, administra zonas, red (WireGuard, MikroTik), planes, inventario, cobranza y usuarios. \\
Técnico (\texttt{tecnico}) & App móvil & Instala en campo (caja NAP, borne, ONU, fotos), atiende averías y retiros, administra su inventario asignado. \\
Vendedor de campo (\texttt{vendedor\_campo}) & App móvil & Capta clientes (datos, GPS, fotos) y registra su jornada y recorrido. \\
Atención al cliente (\texttt{atencion\_cliente}) & Panel web & Caja y mostrador: registra pagos, emite recibos y consulta clientes. \\
Cliente final & App móvil (PIN) y WhatsApp & Consulta su cuenta y pagos; recibe avisos de cobro y corte. \\
\bottomrule
\end{tabularx}
\end{table}

\noindent\textit{Nota de evolución:} el borrador inicial consideraba tres roles
internos (Gerente/Ingeniero, Vendedora, Técnico). Durante la validación se separó el
personal de caja como \emph{Atención al cliente} y se renombró la vendedora como
\emph{Vendedor de campo}; el modelo final tiene exactamente cuatro roles internos
(\texttt{RolUsuario} en \texttt{backend/app/models.py}).

% ======================================================================
\subsection{Requerimientos funcionales}
\label{sec:ce02-rf}

Los requerimientos se agrupan por módulo siguiendo el flujo del negocio: información
pública, venta y validación, instalación en campo, inventario, cobranza, cortes y
comunicación con el cliente.

{\small
\begin{longtable}{L{1.3cm}L{9.6cm}C{1.4cm}}
\caption{Requerimientos funcionales y su estado de implementación.}
\label{tab:ce02-rf}\\
\toprule
\textbf{Código} & \textbf{Descripción} & \textbf{Estado} \\
\midrule
\endfirsthead
\toprule
\textbf{Código} & \textbf{Descripción} & \textbf{Estado} \\
\midrule
\endhead
\midrule \multicolumn{3}{r}{\textit{Continúa en la página siguiente}} \\
\endfoot
\bottomrule
\endlastfoot
\multicolumn{3}{l}{\textbf{A. Portal web público}} \\
RF-01 & Mostrar sin autenticación las zonas con cobertura del servicio. & Impl. \\
RF-02 & Mostrar sin autenticación los planes disponibles (velocidad y precio). & Impl. \\
\multicolumn{3}{l}{\textbf{B. Usuarios y roles}} \\
RF-03 & Inicio de sesión diferenciado para los cuatro roles internos y para el cliente (PIN en la app). & Impl. \\
RF-04 & Restringir cada módulo según el rol autenticado (RBAC). & Impl. \\
\multicolumn{3}{l}{\textbf{C. Registro de ventas (vendedor de campo)}} \\
RF-05 & Registrar DNI/RUC, nombres, teléfono y WhatsApp del cliente, con autocompletado del nombre por consulta de DNI/RUC. & Impl. \\
RF-06 & Capturar automáticamente las coordenadas GPS de la vivienda. & Impl. \\
RF-07 & Adjuntar una o más fotografías de la fachada. & Impl. \\
RF-08 & Registrar dos o más fechas posibles propuestas para la visita técnica. & Pend. \\
\multicolumn{3}{l}{\textbf{D. Validación de ventas (administrador)}} \\
RF-09 & Listar las solicitudes pendientes de aprobación. & Impl. \\
RF-10 & Aprobar o rechazar una solicitud registrando el motivo del rechazo. & Impl. \\
RF-11 & Al aprobar, publicar la orden en el panel del técnico asignado a la zona del cliente. & Impl. \\
\multicolumn{3}{l}{\textbf{E. Coordinación de la visita (técnico)}} \\
RF-12 & El técnico ve únicamente las órdenes aprobadas de sus zonas asignadas. & Impl. \\
RF-13 & Enviar desde el sistema una plantilla de WhatsApp para coordinar la fecha de instalación. & Pend. \\
RF-14 & Registrar la fecha finalmente acordada con el cliente. & Pend. \\
\multicolumn{3}{l}{\textbf{F. Registro de planta externa (técnico)}} \\
RF-15 & Registrar la caja NAP y el borne utilizados en la instalación (la mufa y la corrida de origen se registran en el mapa de red). & Parcial \\
RF-16 & Validar que el borne esté libre antes de asignarlo. & Impl. \\
RF-17 & Vincular automáticamente el borne con el cliente. & Impl. \\
RF-18 & Mostrar el mapa de ocupación de bornes de cada caja (libre / ocupado / dañado). & Impl. \\
\multicolumn{3}{l}{\textbf{G. Registro del equipo instalado (técnico)}} \\
RF-19 & Seleccionar la ONU desde el inventario asignado al técnico. & Impl. \\
RF-20 & Registrar marca y estándar WiFi (WiFi~5 / WiFi~6) del equipo. & Impl. \\
RF-21 & Registrar número de serie y MAC, de forma manual o escaneando el código con la cámara. & Impl. \\
RF-22 & Adjuntar fotografías de evidencia de la instalación. & Impl. \\
RF-23 & Al cerrar la instalación, pasar la ONU a estado \emph{instalado} y descontarla del inventario del técnico. & Impl. \\
\multicolumn{3}{l}{\textbf{H. Inventario de equipos (administrador)}} \\
RF-24 & Registrar nuevas ONU escaneando su código (serie y MAC). & Impl. \\
RF-25 & Asignar ONU a un técnico específico. & Impl. \\
RF-26 & Consultar el stock por técnico y por estado (almacén, asignado, instalado, dañado, perdido). & Impl. \\
RF-27 & Reportar una ONU como dañada indicando el motivo. & Parcial \\
RF-28 & Al reportarse un equipo dañado, notificar al administrador. & Pend. \\
\multicolumn{3}{l}{\textbf{I. Pagos y facturación}} \\
RF-29 & Registrar pagos (efectivo, Yape, transferencia), incluidos adelantos de varios meses con descuento. & Impl. \\
RF-30 & Listar clientes con pagos vencidos o próximos a vencer. & Impl. \\
RF-31 & Generar automáticamente un recibo en PDF por cada cobro. & Impl. \\
RF-32 & Impedir dos pagos con el mismo número de operación. & Impl. \\
\multicolumn{3}{l}{\textbf{J. Cortes de servicio}} \\
RF-33 & Cortar manualmente el servicio de un cliente desde el sistema (address-list en el MikroTik de la zona). & Impl. \\
RF-34 & Ejecutar cortes automáticos diarios para clientes vencidos, respetando días de gracia y compromisos de pago. & Impl. \\
RF-35 & Reactivar automáticamente el servicio al registrar el pago. & Impl. \\
\multicolumn{3}{l}{\textbf{K. Notificaciones por WhatsApp}} \\
RF-36 & Enviar un mensaje de bienvenida al cerrar la instalación (plan, precio, fecha de pago). & Pend. \\
RF-37 & Enviar recordatorios automáticos de pago antes del vencimiento. & Impl. \\
RF-38 & Avisar al cliente antes de ejecutar el corte por falta de pago. & Impl. \\
\multicolumn{3}{l}{\textbf{L. Bot de WhatsApp con IA — validación de pagos}} \\
RF-39 & Recibir la fotografía de un comprobante (Yape/transferencia) enviada por el cliente. & Pend. \\
RF-40 & Solicitar DNI o nombre para ubicar al cliente. & Pend. \\
RF-41 & Extraer con IA (OCR/visión) monto, fecha y número de operación. & Pend. \\
RF-42 & Validar que el comprobante no esté registrado (duplicados). & Pend. \\
RF-43 & Confirmar el pago y enviar el recibo automáticamente. & Pend. \\
\multicolumn{3}{l}{\textbf{M. Autoservicio del cliente}} \\
RF-44 & Consultar la fecha del próximo pago. & Parcial \\
RF-45 & Consultar si tiene cuotas vencidas. & Parcial \\
RF-46 & Consultar el historial de pagos y recibos. & Parcial \\
\multicolumn{3}{l}{\textbf{N. Requerimientos incorporados durante la validación}} \\
RF-47 & Registrar averías (motivo, caja, borne, potencia óptica), asignarlas a un técnico y cerrarlas con diagnóstico y fotos. & Impl. \\
RF-48 & Gestionar órdenes de retiro de equipos (ONU, roseta, patchcord). & Impl. \\
RF-49 & Registrar la jornada del vendedor de campo (inicio/fin, recorrido GPS, pasaje y reintegros). & Impl. \\
RF-50 & Enviar mensajes masivos de WhatsApp por zona y estado, con registro por destinatario. & Impl. \\
RF-51 & Registrar gastos operativos por zona y mostrar ingresos frente a gastos. & Impl. \\
RF-52 & Gestionar clientes satelitales (Starlink) sin zona ni MikroTik. & Impl. \\
RF-53 & Administrar el túnel WireGuard de cada zona (alta/baja de peers, estado del enlace) y una terminal RouterOS restringida. & Impl. \\
RF-54 & Dibujar mufas, corridas de fibra y cajas NAP sobre el mapa de cada zona. & Impl. \\
RF-55 & Sincronizar los planes como perfiles PPPoE en los MikroTik de todas las zonas. & Impl. \\
RF-56 & Mostrar el estado de la sesión PPPoE de un cliente en tiempo real. & Impl. \\
\end{longtable}
}

\noindent Resumen: de 56 requerimientos funcionales, \textbf{41 están implementados},
\textbf{5 parcialmente} y \textbf{10 pendientes}. Los requerimientos RF-44 a RF-46 se
resolvieron mediante la \emph{App del cliente} (consulta de cuenta y pagos con PIN) en
lugar del bot conversacional; el bot con IA (RF-39 a RF-43) se reprograma para el
semestre~2 (ver ADR-06 en el \cref{chap:ce02-arquitectura}).

% ======================================================================
\subsection{Requerimientos no funcionales}
\label{sec:ce02-rnf}

Los requerimientos no funcionales describen \emph{cómo} debe comportarse el sistema, más
allá de lo que hace. Se organizan según las características de calidad de la norma
ISO/IEC~25010 (seguridad, fiabilidad, eficiencia, usabilidad, mantenibilidad y
portabilidad) y se priorizaron a partir de las condiciones reales de WISP: el sistema
maneja datos personales y dinero, por lo que la seguridad y la trazabilidad pesan más; el
técnico trabaja en zonas rurales con señal intermitente, por lo que la usabilidad en el
celular es crítica; y una MYPE con un solo administrador necesita que el sistema se
recupere y se mantenga sin depender de una sola persona.

Cada requerimiento indica cómo se verifica, para que pueda comprobarse con una prueba y no
quede como una intención. La verificación de RNF-10 (recuperación) ya se ejecutó en
laboratorio en el entregable de plataforma de datos (\cref{chap:ce02-datos}).

\begin{table}[H]
\centering
\caption{Requerimientos no funcionales (atributos de calidad ISO/IEC~25010).}
\label{tab:ce02-rnf}
\small
\begin{tabularx}{\textwidth}{L{1.4cm}L{2.3cm}XL{3.3cm}}
\toprule
\textbf{Código} & \textbf{Atributo} & \textbf{Requerimiento} & \textbf{Cómo se verifica} \\
\midrule
RNF-01 & Seguridad & Autenticación con JWT y control de acceso por rol en cada endpoint. & Pruebas de API con roles distintos. \\
RNF-02 & Seguridad & Contraseñas con bcrypt; comunicaciones solo por HTTPS/TLS; freno a la fuerza bruta en el login (8 intentos en 5 min, bloqueo de 15 min). & Revisión de configuración y prueba de bloqueo. \\
RNF-03 & Privacidad & Tratamiento de DNI, ubicación y fotos conforme a la Ley N.\textsuperscript{o}~29733; el token de consulta de DNI vive solo en el servidor. & Auditoría de datos personales. \\
RNF-04 & Disponibilidad & Recordatorios y cortes automáticos operan a diario (08:00 y 09:00, hora de Lima) sin intervención humana. & Registro del planificador. \\
RNF-05 & Usabilidad & La app del técnico funciona en smartphones de gama media y con conectividad intermitente. & Prueba en campo. \\
RNF-06 & Rendimiento & Consultas de borne y estado del cliente en menos de 2~s; el estado de WireGuard se guarda en caché 120~s para no saturar la VPS. & Medición de tiempos de respuesta. \\
RNF-07 & Confiabilidad & Ante duda en un comprobante, el caso pasa a revisión humana (aplica al bot del semestre~2). & Caso de prueba de baja confianza. \\
RNF-08 & Trazabilidad & Todo cambio relevante (activación, corte, pago, cambio de plan u ONU) queda en el historial con usuario y fecha. & Consulta de la tabla \texttt{historial}. \\
RNF-09 & Escalabilidad & Agregar zonas y clientes sin rediseño (zona = registro + peer WireGuard). & Alta de una zona de prueba. \\
RNF-10 & Recuperación & Respaldo diario de la base de datos con retención de 30 días y copia fuera del servidor. & Prueba de restauración. \\
RNF-11 & Mantenibilidad & Backend modular por dominio (19 routers) y esquema versionado con migraciones Alembic. & Revisión de código. \\
RNF-12 & Portabilidad & Todo el sistema se despliega con Docker Compose en cualquier servidor Linux. & Despliegue en un equipo limpio. \\
\bottomrule
\end{tabularx}
\end{table}

% ======================================================================
\subsection{Reglas de negocio}
\label{sec:ce02-rn}

Las reglas de negocio son las condiciones propias de la operación de WISP que el sistema
debe hacer cumplir siempre, sin importar qué usuario opera la pantalla. Se obtuvieron del
modelado de procesos TO-BE (\cref{chap:procesos}) y de la validación con el personal; por
ejemplo, que un borne no puede tener dos clientes, que un técnico solo instala equipos de
su propio inventario o que el aviso de corte se envía \emph{antes} de cortar, para que el
cliente lo reciba mientras todavía tiene internet.

La tabla indica, para cada regla, dónde se garantiza. Las que protegen la integridad de los
datos se refuerzan en la propia base de datos ---restricciones \texttt{UNIQUE}, triggers e
índices probados en el entregable de plataforma de datos---, para que se cumplan aunque una
pantalla tenga un error; las demás las aplica el backend según el rol y el estado del
cliente.

\begin{table}[H]
\centering
\caption{Reglas de negocio y dónde se hacen cumplir.}
\label{tab:ce02-reglas}
\small
\begin{tabularx}{\textwidth}{L{1.2cm}XL{3.4cm}}
\toprule
\textbf{Código} & \textbf{Regla} & \textbf{Dónde se garantiza} \\
\midrule
RN-01 & Un borne de una caja NAP no puede estar asignado a más de un cliente. & Aplicación + base de datos: índice \texttt{uq\_borne\_cliente} y trigger \texttt{trg\_validar\_borne} (\cref{sec:ce02-programacion}). \\
RN-02 & Un técnico solo instala equipos de su propio inventario asignado. & Aplicación. \\
RN-03 & Una instalación no se cierra sin fotos de evidencia y número de serie de la ONU, salvo clientes \emph{migrados}. & Aplicación. \\
RN-04 & Un técnico solo ve y atiende órdenes de sus zonas asignadas. & Aplicación (\texttt{tecnico\_zona}). \\
RN-05 & Una venta debe ser aprobada por el administrador antes de llegar a un técnico. & Aplicación (estado del cliente). \\
RN-06 & No pueden existir dos pagos con el mismo número de operación. & Base de datos (\texttt{UNIQUE}). \\
RN-07 & El corte automático se ejecuta solo si se superan los días de gracia configurados y no hay compromiso de pago vigente. & Planificador. \\
RN-08 & Al registrar el pago de un cliente cortado, el servicio se reactiva automáticamente. & Aplicación + MikroTik; en BD, procedimiento \texttt{sp\_registrar\_pago}. \\
RN-09 & El aviso de corte por WhatsApp se envía \emph{antes} de cortar, para que el cliente lo reciba con internet. & Aplicación. \\
RN-10 & La marca de cliente \emph{migrado} solo la pone quien registra desde el panel, nunca el técnico desde la app. & Aplicación. \\
RN-11 & El usuario y la contraseña PPPoE salen siempre del DNI (nunca del RUC, que es público). & Aplicación. \\
RN-12 & Solo el administrador ve la configuración de red (túnel, router, IP) y ejecuta cortes manuales. & Aplicación (RBAC). \\
\bottomrule
\end{tabularx}
\end{table}

% ======================================================================
\subsection{Restricciones del sistema}
\label{sec:ce02-restricciones}

\begin{itemize}
  \item \textbf{Presupuesto de MYPE:} se prioriza una arquitectura simple y de bajo
  costo operativo (un servidor en la oficina y una VPS pequeña) frente a soluciones
  de alta complejidad.
  \item \textbf{Conectividad de campo:} los técnicos trabajan en zonas rurales con
  cobertura móvil intermitente.
  \item \textbf{Infraestructura heredada:} cada zona ya opera con un MikroTik y
  convenciones propias (nombres de pool PPPoE y de la lista de corte distintos por
  zona); el sistema debe adaptarse sin reconfigurar la red de golpe.
  \item \textbf{Sin IP pública en la oficina:} el sistema se publica mediante un túnel
  de salida (Cloudflare Tunnel) y la red se gestiona a través de una VPS con WireGuard.
  \item \textbf{Marco legal:} Ley N.\textsuperscript{o}~29733 de Protección de Datos
  Personales y su reglamento.
  \item \textbf{Mensajería:} la integración de WhatsApp depende de una sesión
  vinculada y de las políticas de la plataforma (riesgo de bloqueo por envíos masivos).
\end{itemize}

% ======================================================================
\subsection{Historias de usuario y criterios de aceptación}
\label{sec:ce02-hu}

Las historias siguen el formato \emph{Como [rol], quiero [acción], para [beneficio]}
con criterios de aceptación verificables (enfoque \emph{Given/When/Then} resumido).

\newcommand{\historia}[4]{%
\begin{tcolorbox}[enhanced,breakable,colback=white,colframe=upeublue!60,
  boxrule=0.5pt,arc=1pt,left=6pt,right=6pt,top=4pt,bottom=4pt,
  title={\small\textbf{#1} --- #2},coltitle=upeublue,colbacktitle=upeublue!8]
\small #3

\smallskip\textbf{Criterios de aceptación:}
\begin{itemize}[leftmargin=1.2em,itemsep=0pt]
#4
\end{itemize}
\end{tcolorbox}}

\historia{HU-01}{Vendedor de campo · RF-05 a RF-07}
{Como vendedor de campo, quiero registrar los datos, la ubicación y las fotos de un
cliente potencial, para iniciar su contratación sin papeles.}
{\item DNI/RUC, nombre y teléfono son obligatorios; al escribir el DNI el nombre se autocompleta.
 \item Las coordenadas se toman del GPS del dispositivo, no se escriben a mano.
 \item Se adjunta al menos una foto de la fachada.
 \item Al guardar, el cliente queda en estado \texttt{pendiente\_aprobacion}.}

\historia{HU-02}{Administrador · RF-09 a RF-11}
{Como administrador, quiero aprobar o rechazar las ventas registradas, para asegurar
que los datos son correctos antes de enviar un técnico.}
{\item El panel lista los clientes en \texttt{pendiente\_aprobacion}.
 \item Aprobar cambia el estado a \texttt{aprobado}; rechazar exige un motivo.
 \item Al aprobar, la orden aparece en la app del técnico asignado a esa zona.}

\historia{HU-03}{Técnico · RF-12}
{Como técnico, quiero ver solo las instalaciones de mis zonas, para saber qué me
corresponde sin depender de mensajes externos.}
{\item El listado filtra por las zonas asignadas en \texttt{tecnico\_zona}.
 \item Cada orden muestra los datos cargados por el vendedor y su ubicación en el mapa.}

\historia{HU-04}{Técnico · RF-15 a RF-18}
{Como técnico, quiero registrar la caja NAP y el borne que uso, para dejar constancia
exacta de la conexión física del cliente.}
{\item El sistema muestra los bornes de la caja diferenciando libres, ocupados y dañados.
 \item No se puede elegir un borne ocupado.
 \item Al confirmar, el borne queda vinculado al cliente y en estado \texttt{ocupado}.}

\historia{HU-05}{Técnico · RF-19 a RF-23}
{Como técnico, quiero registrar la ONU que dejo instalada desde mi inventario, para que
el sistema controle el stock y la trazabilidad por número de serie.}
{\item Solo aparecen ONU con estado \texttt{asignado\_tecnico} a ese técnico.
 \item La serie/MAC se puede escanear con la cámara.
 \item Al cerrar la instalación, la ONU pasa a \texttt{instalado} y el cliente a \texttt{activo}.}

\historia{HU-06}{Técnico · RF-47}
{Como técnico, quiero atender averías con diagnóstico, potencia óptica y fotos, para
resolverlas y dejar evidencia del trabajo.}
{\item La avería pasa por \texttt{pendiente} $\rightarrow$ \texttt{asignada} $\rightarrow$ \texttt{en\_proceso} $\rightarrow$ \texttt{resuelta}.
 \item Si se reemplaza la ONU, se registran la retirada y la nueva.}

\historia{HU-07}{Administrador · RF-24 a RF-26}
{Como administrador, quiero registrar y asignar ONU escaneando su código, para
controlar el inventario sin digitar cada número de serie.}
{\item Al escanear, el sistema completa la serie y permite corregirla si la lectura falla.
 \item El stock se consulta por técnico y por estado.}

\historia{HU-08}{Atención al cliente · RF-29 a RF-32}
{Como personal de atención, quiero registrar pagos y emitir el recibo, para cobrar en
mostrador de forma rápida y sin errores.}
{\item Se puede registrar un adelanto de varios meses con un solo recibo.
 \item Un número de operación repetido se rechaza.
 \item Si el cliente estaba cortado, el servicio se reactiva al registrar el pago.}

\historia{HU-09}{Administrador · RF-33 a RF-35}
{Como administrador, quiero que los cortes por mora sean automáticos y poder cortar
manualmente, para gestionar la cobranza sin entrar a cada router.}
{\item Cada día a las 09:00 se cortan los clientes vencidos más allá de los días de gracia.
 \item Se respetan los compromisos de pago vigentes.
 \item El corte agrega la IP del cliente a la lista de corte configurada para su zona.}

\historia{HU-10}{Cliente · RF-37, RF-38}
{Como cliente, quiero recibir avisos antes del vencimiento y antes del corte, para
pagar a tiempo y no quedarme sin servicio por olvido.}
{\item El recordatorio se envía unos días antes del vencimiento (08:00).
 \item El aviso de corte llega antes de que se ejecute el corte.}

\historia{HU-11}{Cliente · RF-44 a RF-46}
{Como cliente, quiero consultar en la app cuándo vence mi servicio y mis pagos, para
resolver mis dudas sin llamar a la empresa.}
{\item El cliente entra con su PIN.
 \item Ve su plan, el estado de su servicio, su fecha de vencimiento y el historial de pagos.}

\historia{HU-12}{Vendedor de campo · RF-49}
{Como vendedor de campo, quiero registrar mi jornada y recorrido, para que la empresa
vea mi trabajo y me reintegre el pasaje.}
{\item La jornada guarda el punto de inicio, los puntos del recorrido y el de fin.
 \item El administrador ve el recorrido en el mapa y marca el reintegro.}

% ======================================================================
\subsection{Validación con los interesados}
\label{sec:ce02-validacion}

La especificación y el prototipo (\cref{chap:ce02-prototipo}) se revisaron con los
perfiles de usuario de WISP (gerencia/ingeniería, ventas, técnico de campo y atención).
Los cambios acordados ya forman parte de este documento:

\begin{itemize}
  \item La ubicación se captura por GPS y no como texto, para evitar errores de tipeo (RF-06).
  \item El técnico pidió ver la caja NAP como un mapa visual de bornes y no como lista (RF-18).
  \item El administrador pidió que, al escanear un equipo, el dato leído se muestre
  de inmediato para corregirlo si la lectura falla (RF-24).
  \item Se agregaron los módulos de averías, retiros, jornadas, mensajes masivos,
  gastos y clientes satelitales (RF-47 a RF-52), que el negocio necesitaba y no
  estaban en el borrador.
  \item El cliente prefirió una app con su cuenta en lugar de esperar al bot; el bot
  con IA queda para una iteración posterior.
\end{itemize}

\begin{table}[htbp]
\centering
\caption{Acta de validación de requerimientos (Anexo de la evidencia CE0211).}
\label{tab:ce02-acta}
\small
\begin{tabularx}{\textwidth}{L{3.6cm}X}
\toprule
\textbf{Campo} & \textbf{Detalle} \\
\midrule
Proyecto & Sistema de Gestión Integral para ISP de Fibra Óptica Multi-Zona (WISP) \\
Objetivo de la sesión & Validar la especificación de requerimientos y el prototipo navegable. \\
Fecha y lugar & \rule{4cm}{0.4pt} \quad Oficina central de WISP, Juliaca \\
Participantes & Gerente/Ingeniero, vendedor de campo, técnico de campo, personal de atención, equipo del proyecto \\
Resultado & Requerimientos aprobados con las observaciones listadas arriba. \\
\bottomrule
\end{tabularx}

\medskip
\begin{tabularx}{\textwidth}{XXX}
\rule{4.2cm}{0.4pt} & \rule{4.2cm}{0.4pt} & \rule{4.2cm}{0.4pt} \\
Gerente / Ingeniero & Técnico de campo & Vendedor de campo \\
\end{tabularx}
\end{table}

\begin{upeubox}{Pendiente antes de la entrega}
El acta debe completarse con la fecha real de la sesión y las firmas de los
participantes. Un acta sin firmas no se considera evidencia de validación.
\end{upeubox}

% ======================================================================
\subsection{Matriz de trazabilidad de requerimientos}
\label{sec:ce02-trazabilidad}

{\small
\begin{longtable}{L{2.5cm}L{1.4cm}L{4.6cm}L{4.4cm}}
\caption{Trazabilidad requerimiento $\rightarrow$ historia $\rightarrow$ componente $\rightarrow$ diseño/datos.}
\label{tab:ce02-trazabilidad}\\
\toprule
\textbf{Requerimientos} & \textbf{HU} & \textbf{Componente (router / pantalla)} & \textbf{Diseño y datos} \\
\midrule
\endfirsthead
\toprule
\textbf{Requerimientos} & \textbf{HU} & \textbf{Componente (router / pantalla)} & \textbf{Diseño y datos} \\
\midrule
\endhead
\bottomrule
\endlastfoot
RF-01, RF-02 & --- & Sitio público (\texttt{web/}), \texttt{zonas.py} (\texttt{/publicas}) & Tablas \texttt{zonas}, \texttt{planes} \\
RF-03, RF-04 & --- & \texttt{auth.py}, \texttt{cliente\_app.py} & Tabla \texttt{usuarios}; \texttt{clientes.pin\_app} \\
RF-05 a RF-07 & HU-01 & \texttt{clientes.py}, \texttt{consultas.py}; app \texttt{(vendedor)/nuevo-cliente} & Secuencia de venta; tabla \texttt{clientes} \\
RF-09 a RF-11 & HU-02 & \texttt{clientes.py}; panel \texttt{AprobarClienteDialog} & Estados del cliente \\
RF-12 & HU-03 & \texttt{clientes.py}; app \texttt{(tecnico)/instalaciones} & Tabla \texttt{tecnico\_zona} \\
RF-15 a RF-18 & HU-04 & \texttt{cajas\_nap.py}, \texttt{mufas.py}; app \texttt{(tecnico)/instalar} & Secuencia de instalación; tablas \texttt{cajas\_nap}, \texttt{bornes} \\
RF-19 a RF-26 & HU-05, HU-07 & \texttt{onus.py}; app \texttt{escanear-equipo}, \texttt{inventario} & Estados de la ONU; tabla \texttt{onus} \\
RF-29 a RF-32 & HU-08 & \texttt{pagos.py}, \texttt{recibos.py} & Tabla \texttt{pagos} (\texttt{UNIQUE numero\_operacion}) \\
RF-33 a RF-35 & HU-09 & \texttt{clientes.py}, \texttt{scheduler.py} & Secuencia de corte; tabla \texttt{historial} \\
RF-37, RF-38, RF-50 & HU-10 & \texttt{whatsapp.py}, servicio \texttt{whatsapp/} & Tablas \texttt{mensajes\_masivos*} \\
RF-44 a RF-46 & HU-11 & \texttt{cliente\_app.py}; app \texttt{(cliente)/} & Tablas \texttt{clientes}, \texttt{pagos} \\
RF-47, RF-48 & HU-06 & \texttt{averias.py}, \texttt{ordenes\_retiro.py}, \texttt{evidencias.py} & Tablas \texttt{averias}, \texttt{ordenes\_retiro}, \texttt{evidencias} \\
RF-49 & HU-12 & \texttt{jornadas.py}; app \texttt{(vendedor)/jornada} & Tablas \texttt{jornadas\_vendedor}, \texttt{ubicaciones\_vendedor} \\
RF-51 & --- & \texttt{gastos.py}, \texttt{dashboard.py} & Tabla \texttt{gastos} \\
RF-53, RF-55, RF-56 & --- & \texttt{zonas.py}, \texttt{planes.py} & Despliegue (WireGuard); tabla \texttt{zonas} \\
\end{longtable}
}
